% chapters/07/evaluasi.tex
\bagian{Asesmen dan Instrumen Penilaian}
\label{sec:asesmen-tabrakan}

\begin{asesmenbox}
\textbf{Kedudukan asesmen bab ini.} Bab ini menuntaskan \SUBCPMK{1.4} sekaligus
menutup \CPMK{1}. Ketiga indikatornya diperiksa melalui penilaian artefak
program (Tugas Praktikum 05) dan soal konsep yang menuntut analisis.

\vspace{0.4em}
Materi Bab~VII termasuk bahan \textbf{Ujian Tengah Semester} pada minggu
kesepuluh, bersama Bab~III sampai Bab~XI untuk \CPMK{1} dan \CPMK{2}.
\end{asesmenbox}

\subsection*{Instrumen 1: Rubrik Praktikum 05}

\begin{table}[htbp]
\centering
\caption{Rubrik penilaian Praktikum 05 beserta indikator yang diukur}
\label{tab:rubrik-praktikum-05}
{\setlength{\tabcolsep}{4pt}\footnotesize%
\begin{tabularx}{\textwidth}{@{}L{3.2cm} C{1.3cm} C{2.1cm} Y@{}}
\toprule
\headrow \textbf{Kriteria} & \textbf{Bobot} & \textbf{Indikator} & \textbf{Deskripsi capaian} \\
\midrule
Penyelesaian tabrakan per sumbu & 35\% & \ind{b} &
Pemain dapat mendarat di atas bongkahan dan terhalang oleh sisinya; penghitung
tumpang tindih bernilai nol pada seluruh laju bingkai yang diuji. \\
\addlinespace
Struktur kelas dan wadah & 30\% & \ind{a} &
Benda disusun sebagai kelas turunan \code{Sprite}, dikelola dengan
\code{Group}, dan perintah pemeriksaan tabrakan dipilih sesuai keperluan. \\
\addlinespace
Analisis kesalahan tabrakan & 35\% & \ind{c} &
Catatan praktikum memuat ketiga gejala hasil percobaan perusakan, penyebabnya,
dan perbaikan yang dilakukan, disertai penjelasan pada kecepatan tinggi. \\
\bottomrule
\end{tabularx}}
\end{table}

\subsection*{Instrumen 2: Soal Analisis}

Soal berikut menuntut analisis, bukan sekadar ingatan. Kunci jawaban beserta
indikator yang diukur disajikan pada Tabel~\ref{tab:kunci-konsep-ch07}.

\begin{enumerate}
  \item Sebuah benda bergerak 1200 piksel per detik pada laju 60 bingkai per
        detik, dan bongkahan yang dilaluinya tebalnya 15 piksel. Jelaskan apa
        yang terjadi, sertakan perhitungannya, dan sebutkan dua cara
        perbaikan.
  \item Seorang pemain berjalan ke kanan dan menabrak sisi kiri bongkahan,
        tetapi tiba-tiba muncul di atas bongkahan itu. Jelaskan penyebabnya
        dan bagaimana Anda menemukannya.
  \item Pemain bergetar halus ketika berdiri di atas bongkahan. Jelaskan
        penyebabnya dan perbaikannya.
  \item Sebuah permainan memuat benda berbentuk huruf melengkung. Metode
        pemeriksaan tabrakan apa yang Anda pilih, dan mengapa?
  \item Jelaskan bagaimana Anda memastikan bahwa penanganan tabrakan pada
        program Anda memang benar, tanpa harus memainkannya berulang kali.
\end{enumerate}

\begin{table}[htbp]
\centering
\caption{Kunci soal analisis Bab VII beserta indikator yang diukur}
\label{tab:kunci-konsep-ch07}
{\setlength{\tabcolsep}{4pt}\footnotesize%
\begin{tabularx}{\textwidth}{@{}C{0.8cm} C{2.1cm} Y@{}}
\toprule
\headrow \textbf{No} & \textbf{Indikator} & \textbf{Pokok jawaban} \\
\midrule
1 & \ind{c} &
Perpindahan setiap bingkai $1200 \div 60 = 20$ piksel, lebih panjang daripada
tebal bongkahan 15 piksel, sehingga benda melompati bongkahan tanpa pernah
beririsan pada saat pemeriksaan. Perbaikan: batasi panjang langkah, bagi
perpindahan menjadi beberapa langkah, atau tebalkan bongkahan. \\
\addlinespace
2 & \ind{c} &
Kedua sumbu diperiksa sekaligus sehingga program selalu menyelesaikan pada
sumbu tegak, sebab kecepatan tegaknya bernilai positif. Penelusuran dilakukan
dengan menampilkan nilai \code{rect} sebelum dan sesudah pembaruan pada
bingkai yang bermasalah, lalu memeriksa sumbu yang diperbaiki. \\
\addlinespace
3 & \ind{b} &
Kecepatan tegak tidak dinolkan ketika tabrakan diselesaikan, sehingga benda
terus tertekan ke dalam bongkahan dan posisinya diperbaiki pada setiap
bingkai. Perbaikan: nolkan kecepatan pada sumbu yang diselesaikan. \\
\addlinespace
4 & \ind{a} &
Persegi pembatas tidak memadai karena bentuknya jauh lebih kecil daripada
persegi pembatasnya, sehingga sering terjadi tabrakan palsu. Metode yang
tepat adalah \istilah{mask}, dengan persegi pembatas dipakai lebih dahulu
sebagai penyaring agar \istilah{mask} hanya dijalankan pada benda yang
berdekatan. \\
\addlinespace
5 & \ind{c} &
Dengan menambahkan pemeriksa mandiri di dalam program, misalnya penghitung
bongkahan yang bertumpang tindih dengan pemain, lalu menjalankan program
tanpa sentuhan manusia pada beberapa laju bingkai dan memastikan nilainya
selalu nol. \\
\bottomrule
\end{tabularx}}
\end{table}
